\begin{table*}[t]\centering\caption{Main results (five seeds, mean $\pm$ std, relative $L_2$; lower is better): test error at the training grid (test 128), training error (train 128), zero-shot test errors at 256 and 512 points with the native rules of Section~\ref{sec:method} (test 256, test 512), the 256-point test error under the resampled rule (rs256), parameter count, and wall-clock training time.}\label{tab:main}
\begin{tabular}{lcccccrr}
\toprule
System & test 128 & train 128 & test 256 & test 512 & rs256 & params & time (s) \\
\midrule
FNO & 0.0049 $\pm$ 0.0004 & 0.0031 $\pm$ 0.0001 & 0.0049 $\pm$ 0.0004 & 0.0049 $\pm$ 0.0004 & 0.0049 $\pm$ 0.0004 & 139745 & 68.5 $\pm$ 10.5 \\
DeepONet & 0.2192 $\pm$ 0.0168 & 0.1988 $\pm$ 0.0166 & 0.2195 $\pm$ 0.0168 & 0.2196 $\pm$ 0.0168 & 0.2192 $\pm$ 0.0168 & 66305 & 2.9 $\pm$ 0.2 \\
MLP & 0.0834 $\pm$ 0.0069 & 0.0236 $\pm$ 0.0007 & 0.0834 $\pm$ 0.0069 & 0.0834 $\pm$ 0.0069 & 0.0834 $\pm$ 0.0069 & 394368 & 5.5 $\pm$ 0.5 \\
CNN & 0.2229 $\pm$ 0.0078 & 0.2103 $\pm$ 0.0055 & 0.5679 $\pm$ 0.0276 & 0.7270 $\pm$ 0.0371 & 0.2229 $\pm$ 0.0078 & 37601 & 166.7 $\pm$ 33.4 \\
\bottomrule
\end{tabular}
\end{table*}

\begin{table}[t]\centering\caption{Zero-shot error ratio (mean over seeds of fine-grid over 128-point test error).}\label{tab:ratio}
\begin{tabular}{lrr}
\toprule
System & 256/128 & 512/128 \\
\midrule
FNO & 1.00 & 1.00 \\
DeepONet & 1.00 & 1.00 \\
MLP & 1.00 & 1.00 \\
CNN & 2.55 & 3.27 \\
\bottomrule
\end{tabular}
\end{table}

\begin{table}[t]\centering\caption{H1: FNO versus each baseline at 128 points (Welch test over five seeds).}\label{tab:h1}
\begin{tabular}{lrrrr}
\toprule
Baseline & mean & FNO gain & Welch $p$ & Cohen $d$ \\
\midrule
CNN & 0.2229 & 97.8\% & 3.7e-07 & 39.7 \\
DeepONet & 0.2192 & 97.8\% & 9.0e-06 & 18.0 \\
MLP & 0.0834 & 94.1\% & 1.4e-05 & 16.0 \\
\bottomrule
\end{tabular}
\end{table}

\begin{figure}[t]\centering\includegraphics[width=\columnwidth]{resolution_transfer.pdf}
\caption{Test relative $L_2$ at the training grid and zero-shot grids (mean, error bars one std over five seeds, log scale).}\label{fig:res}\end{figure}

\textbf{H1 (accuracy at the training grid): supported.} The FNO error is far below all three baselines (Table~\ref{tab:main}); the gain over the best baseline (MLP) is far above the 10\% threshold and all three Welch tests pass at $\alpha=0.05$ (Table~\ref{tab:h1}). DeepONet and the CNN are roughly equal to each other and more than an order of magnitude worse than the FNO. What this does not show: that DeepONet or the CNN are intrinsically that weak. Their training errors are also high (Table~\ref{tab:main}), so both are under-fit under the 100-epoch budget. For DeepONet, the training-length sweep (Section~\ref{sec:ablations}) lowers the test error to 0.1314 at 500 epochs and 0.1276 at 2000, still far above the FNO. This does not reproduce the comparable performance of DeepONet and FNO that Lu et al.~\cite{lu2021comprehensive} report for simple settings: our DeepONet is the plain architecture with one tuned hyperparameter, we compare against no published error value, and its numbers are a lower bound on DeepONet quality. For the CNN, a structural limit may add to the under-fitting: its receptive field covers 49 of the 128 grid points. We ran no CNN with a larger receptive field or longer training, so we cannot separate the two causes. The MLP, by contrast, fits its training set much better than its test set, i.e.\ it overfits 400 samples.

\textbf{H2 (zero-shot finer grid): supported as stated, with a caveat that changes its meaning.} The FNO ratio is 1.00 at both finer grids and the CNN ratio exceeds 1.5 (Table~\ref{tab:ratio}, Fig.~\ref{fig:res}). The FNO error at 256 and 512 points equals its 128-point error to the digits shown, which is what one expects here: the initial conditions are band-limited to 63 modes, the solver resolves the solution on all three grids, and a spectral model sees the same function at each resolution. The test is therefore an easy case of resolution invariance and says nothing about functions with content above the training Nyquist frequency. The CNN degrades when it is applied natively. A likely cause is that its kernel spans a fixed number of grid points, i.e.\ a smaller physical width on a finer grid; we did not isolate this with a separate run. Under the resampled rule the CNN's 256-point error (rs256 in Table~\ref{tab:main}) matches its 128-point error, so the degradation comes from native application to the finer input and not from the fine-grid targets. The hypothesis also predicted that the MLP and DeepONet would not transfer; that part is \emph{refuted}: their ratio is 1.00, but only because under our rule they receive the 128 stride-subsampled input values and so cannot see the finer grid at all. Their zero-shot accuracy is their 128-point accuracy, not evidence of resolution invariance.

\textbf{Failure cases.} The CNN applied natively on finer grids is the clear failure: its 256- and 512-point errors are far above those of the MLP and, at 512, more than three times its own training-grid error (Table~\ref{tab:ratio}). DeepONet and the CNN also fail at the task itself under this budget (errors above 0.2 at the training grid, with training error of the same size). Finally, a finer grid does not help the FNO either: its error is the same on all three grids to the digits shown, so on this data resolution is not what limits it.
